\subsection{问题二的建模与求解：双角色散消干扰}

两条碳化硅光谱来自同一块晶圆，逐行共享$7469$个波数点，却有$2.14$个百分点的反射率均方根差；这意味着厚度应当共享，基线、振幅和相位却不能把两角硬压成一条曲线。双角色散消干扰据此把共同的条纹位置交给厚度与色散，把角度特有的慢变响应留在各自的有限参数中，而不是将两角原谱先平均。
% src:交接/数据档案.json:跨文件关系.同片分组
% src:交接/数据档案.json:跨文件关系.两两全量核验

全量同片拟合给出$10.80\ \mu\mathrm{m}$的拟合厚度；五折连续留段显示厚度分支明显漂移，因此谱形解释与厚度辨识并列呈现。正式留段标准化均方根误差为$1.21$，相对二次趋势基线改善$16.2\%$，说明条纹项增加了模型对留出反射率的解释力。
% src:求解/问题2/结果/厚度结果.json:同片最佳厚度_微米
% src:求解/问题2/结果/验证.json:正式主方法分数_连续留段标准化均方根误差;正式主方法相对正式二次趋势改善_百分比

这条主线把“共享厚度”与“角度响应”分成两个层次，下面依次给出分块准备、共同厚度投影模型及其结果边界。
